% Auto-generated by scripts/latex_synthesis.py -- do not hand-edit; regenerate from the
% source synthesis/<slug>.md instead, since edits here will be overwritten on the next run.
%
% Thesis-chapter style: report class (numeric \cite by default, no natbib/biblatex needed),
% standard margins, no title page/TOC -- this is meant to be read as a single chapter, not a
% standalone thesis. See scripts/latex_synthesis.py's CITATION REPLACEMENT STRATEGY comment
% for how the "Author et al. (Year)" text was turned into \cite{} calls below.
\documentclass[12pt]{report}
\usepackage[margin=1in]{geometry}
\usepackage[utf8]{inputenc}
\usepackage[T1]{fontenc}

\begin{document}

\chapter{how is transfer learning / pretraining used across these papers}

% Cross-paper synthesis generated by scripts/synthesize.py.

Across these five papers, transfer learning and pretraining emerge as a central technique specifically in the language and vision papers, while the interpretability and translation-architecture papers use pretrained models only as objects of analysis or as a stated limitation to be addressed. \cite{1810.04805} makes pretraining the paper's core contribution, distinguishing between the feature-based approach exemplified by ELMo, which incorporates pretrained representations as additional task-specific features, and the fine-tuning approach exemplified by OpenAI GPT, in which nearly all parameters are initialized from a pretrained model and then adapted end to end on the downstream task; BERT extends this second paradigm by introducing a masked language model objective that enables bidirectional pretraining, which the authors argue is what allows a single pretrained architecture to transfer effectively to both sentence-level and token-level tasks with minimal architecture change \cite{1810.04805}. \cite{1810.04805} explicitly situates this within a longer history of transfer learning in NLP, citing prior work in unsupervised feature-based pretraining of word embeddings as well as supervised transfer from tasks such as natural language inference and machine translation, and it further draws an analogy to computer vision, noting that fine-tuning models pretrained on ImageNet had already established the value of large-scale pretrained representations in that domain \cite{1810.04805} — a parallel made concrete by \cite{1512.03385}, whose deep residual networks, pretrained on ImageNet, are shown to transfer well to other recognition tasks such as COCO object detection, yielding substantial accuracy gains from representations learned at very large depth \cite{1512.03385}.

By contrast, \cite{10-1371-journal-pone-0130140} treats a pretrained classifier not as a technique to be extended but as a fixed artifact to be explained: the paper evaluates its layer-wise relevance propagation method on an already pretrained ImageNet model from the Caffe package, explicitly noting that its contribution is not about classifier training and that the proposed methods are built on top of, and applicable to, an already pretrained classifier \cite{10-1371-journal-pone-0130140}. \cite{1409.0473} and \cite{1706.03762}, meanwhile, address transfer learning only tangentially: \cite{1409.0473} trains its RNN-based encoder–decoder models from scratch on parallel corpora, though it briefly notes that pretraining an encoder on a larger monolingual corpus might be possible and beneficial, without pursuing this further, while \cite{1706.03762} similarly trains the Transformer from scratch for translation and parsing rather than employing any pretraining step — a gap later exploited by \cite{1810.04805}, which adopts the Transformer encoder architecture from \cite{1706.03762} as the backbone for its own pretraining framework.

\bibliographystyle{plain}
\bibliography{how-is-transfer-learning-pretraining-used-across-these-paper}

\end{document}
